% This document was generated by an LLM.

\documentclass{essay}

\title{Three Theorems and One Axiom}
\author{An essay for the reader leaving calculus for analysis}
\date{}

\begin{document}
\maketitle

\section{What changes}

In calculus you were told that a continuous function on a closed interval attains a maximum. You probably nodded. The picture is compelling: a curve drawn over a finite stretch of the axis has to have a highest point, because you can see it. The eye finds the peak before the mind asks for a reason.

Analysis begins at the moment you notice that the picture is not an argument, and---worse---that the picture is a \emph{claim about the real numbers} that you have never verified. Consider the same statement over the rationals. Let $f(x) = x^{2}$ on the rational points of $[0,2]$, but delete from the domain every rational $x$ with $x^{2} > 2$. What remains is a set of rationals with no largest element, and $f$ on it is perfectly continuous in the $\varepsilon$--$\delta$ sense, bounded, and attains no maximum. Nothing in the definition of continuity forbids this. The failure is not in $f$. It is in the number system, which has a hole where $\sqrt{2}$ should be.

That is the entire content of what follows. The three theorems of this essay---the Intermediate Value Theorem, the Extreme Value Theorem, the Mean Value Theorem---are not facts about functions. They are facts about $\mathbb{R}$, wearing functions as clothing. Each one is false over $\mathbb{Q}$. Each one is proved by locating a hole where there is none, and deriving a contradiction from the assumption that the hole exists.

A note on the collective name. These are sometimes called ``the three basic theorems of calculus,'' which is misleading; the Fundamental Theorem of Calculus is a different animal, about integration, and does not sit in their logical chain. What actually unites the three is the axiom below.

\section{The axiom, and the two definitions}

Everything rests on one sentence.

\begin{axiom}[Completeness]
Every nonempty set $S \subseteq \mathbb{R}$ that is bounded above has a least upper bound $\sup S$, and $\sup S \in \mathbb{R}$.
\end{axiom}

This is not a theorem. It is an axiom, or in a constructed model of $\mathbb{R}$ it is the property the construction is designed to deliver. It is the only thing distinguishing $\mathbb{R}$ from $\mathbb{Q}$ as an ordered field, and it is the only ingredient in this essay you cannot see on a graph.

Two definitions carry it into the theorems.

\begin{definition}[Continuity]
$f : D \to \mathbb{R}$ is continuous at $c \in D$ if for every $\varepsilon > 0$ there exists $\delta > 0$ such that $x \in D$ and $|x - c| < \delta$ imply $|f(x) - f(c)| < \varepsilon$.
\end{definition}

Observe what this definition does \emph{not} say. It says nothing about drawing without lifting the pen. It is a statement about the order of two quantifiers: the $\varepsilon$ is chosen by an adversary, the $\delta$ by you, and you must answer every challenge. Note also that quantifying $x$ over $D$ makes endpoint continuity automatically one-sided; you do not need a separate definition for it.

\begin{definition}[Differentiability]
$f$ is differentiable at an interior point $c$ of its domain if
\[
f'(c) = \lim_{x \to c} \frac{f(x) - f(c)}{x - c}
\]
exists in $\mathbb{R}$.
\end{definition}

``Interior'' is doing real work here. The quotient must be approachable from both sides.

From continuity alone, one small lemma. It is the workhorse of Section~3, and it is worth seeing how little it costs.

\begin{lemma}[Local sign preservation]\label{lem:sign}
If $f$ is continuous at $c$ and $f(c) \neq 0$, then $f$ has the sign of $f(c)$ throughout some neighbourhood of $c$.
\end{lemma}

\begin{proof}
Take the adversary's $\varepsilon$ to be $|f(c)|$, which is positive by hypothesis. Continuity hands you a $\delta$ with $|f(x) - f(c)| < |f(c)|$ for all $x$ in the domain with $|x-c| < \delta$. If $f(c) > 0$ this forces $f(x) > f(c) - |f(c)| = 0$; if $f(c) < 0$ it forces $f(x) < f(c) + |f(c)| = 0$.
\end{proof}

That is a complete proof, and it is three lines long. The technique---choose $\varepsilon$ to be a specific quantity built from the data, not an anonymous small number---is the single most transferable habit in this essay.

\section{The Intermediate Value Theorem}

\begin{theorem}[Intermediate Value Theorem]\label{thm:ivt}
Let $f : [a,b] \to \mathbb{R}$ be continuous with $f(a) < \gamma < f(b)$. Then $f(c) = \gamma$ for some $c \in (a,b)$.
\end{theorem}

The calculus argument is: the graph starts below the line $y = \gamma$ and ends above it, so it crosses. The analysis argument replaces ``crosses'' with a construction. Where, exactly, does it cross? At the \emph{last moment before it is finished being below}. That phrase is a supremum.

\begin{proof}
Put
\[
S = \{\, x \in [a,b] : f(x) < \gamma \,\}.
\]
Then $a \in S$, so $S$ is nonempty, and $b$ bounds it above. Completeness gives $c = \sup S \in [a,b]$. We show both alternatives to $f(c) = \gamma$ are impossible.

Suppose $f(c) < \gamma$. Then $c \neq b$. Apply Lemma~\ref{lem:sign} to $g = f - \gamma$, which is continuous at $c$ with $g(c) < 0$: there is $\delta > 0$ on which $f$ stays below $\gamma$. Since $c < b$, the point $x_{1} = \min(b, c + \delta/2)$ lies in $S$ and exceeds $c$---contradicting that $c$ is an \emph{upper bound} for $S$.

Suppose $f(c) > \gamma$. Then $c \neq a$. Lemma~\ref{lem:sign} gives $\delta > 0$ on which $f$ stays above $\gamma$; shrink $\delta$ so that $c - \delta \geq a$. Then $(c - \delta, c]$ misses $S$ entirely, so every element of $S$ is at most $c - \delta$, making $c - \delta$ an upper bound smaller than $c$---contradicting that $c$ is the \emph{least} upper bound.

By trichotomy, $f(c) = \gamma$; and $c$ is interior since $f(a) \neq \gamma \neq f(b)$. For the descending case apply the above to $-f$ and $-\gamma$.
\end{proof}

Look at the shape of that argument. The supremum has two defining properties---upper bound, and least such---and the two failure cases contradict exactly one each. This is not a coincidence, and it is not a trick particular to the Intermediate Value Theorem. When a proof invokes $\sup$, expect both halves of the definition to be consumed before it ends. If only one is used, the hypothesis was probably stronger than necessary.

\section{The Extreme Value Theorem}

Boundedness comes first, and it is a genuine theorem, not a preliminary. On $(0,1]$, the continuous function $1/x$ is unbounded. Compactness of the domain is not decoration.

\begin{lemma}[Boundedness]\label{lem:bdd}
A continuous $f : [a,b] \to \mathbb{R}$ is bounded.
\end{lemma}

\begin{proof}
Let $S = \{\, x \in [a,b] : f \text{ is bounded on } [a,x] \,\}$. Then $a \in S$ and $b$ bounds $S$; set $c = \sup S$.

Continuity at $c$ with $\varepsilon = 1$ gives $\delta > 0$ such that $|f| \leq |f(c)| + 1$ on $I = (c - \delta, c + \delta) \cap [a,b]$. Because $c - \delta$ is less than the least upper bound, some $x_{0} \in S$ satisfies $x_{0} > c - \delta$; let $M_{0}$ bound $|f|$ on $[a, x_{0}]$. Now for any $y \in [a,b]$ with $y < c + \delta$ we have $[a,y] \subseteq [a,x_{0}] \cup I$, so $|f| \leq \max(M_{0}, |f(c)| + 1)$ there, and $y \in S$.

Hence $c \in S$. And if $c < b$, then $\min(b, c + \delta/2)$ is a point of $S$ beyond $c$, which is absurd. So $c = b \in S$.
\end{proof}

The set $S$ here is ``the points up to which the theorem is already true.'' Defining such a set and proving its supremum is both attained and maximal is the discrete skeleton of an induction, run along a continuum. You will meet this pattern again under the name of connectedness arguments.

\begin{theorem}[Extreme Value Theorem]\label{thm:evt}
A continuous $f : [a,b] \to \mathbb{R}$ attains a maximum and a minimum: there are $p, q \in [a,b]$ with $f(p) \leq f(x) \leq f(q)$ for all $x \in [a,b]$.
\end{theorem}

\begin{proof}
By Lemma~\ref{lem:bdd} the image is bounded and nonempty, so $M = \sup f([a,b])$ exists. Suppose, for contradiction, that $f(x) < M$ everywhere. Then
\[
g(x) = \frac{1}{M - f(x)}
\]
has a continuous, nowhere-vanishing denominator, so $g$ is continuous on $[a,b]$, hence bounded by Lemma~\ref{lem:bdd}, say $g \leq K$ with $K > 0$. Unwinding: $M - f(x) \geq 1/K$, so $f(x) \leq M - 1/K$ for every $x$. That exhibits an upper bound for the image strictly below $M$---impossible. So $f(q) = M$ for some $q$. Apply the result to $-f$ for the minimum.
\end{proof}

That $g$ is a magic trick, and you should be suspicious of magic tricks. The honest reading: the assumption ``$f$ never reaches its supremum'' is precisely the assumption that $M - f$ is positive but arbitrarily small, and dividing is how you convert \emph{arbitrarily small} into \emph{arbitrarily large}, where you already have a theorem. Every construction of this kind is a translation between two forms of a single failure.

\section{Rolle, and the Mean Value Theorem}

The Mean Value Theorem is where derivatives enter, and it is the only theorem here that depends on another. Its chain is: the Extreme Value Theorem gives an extremum; the extremum is interior; interior extrema have vanishing derivative.

The middle link is Fermat's lemma, and it is the cleanest example in elementary analysis of a limit inheriting an inequality.

\begin{lemma}[Interior extremum]\label{lem:fermat}
If $c \in (a,b)$ is a local extremum of $f$ and $f'(c)$ exists, then $f'(c) = 0$.
\end{lemma}

\begin{proof}
Say $c$ is a local maximum, so $f(x) \leq f(c)$ on some $(c - \eta, c + \eta) \subseteq (a,b)$. For $x$ slightly to the right of $c$, the difference quotient has numerator $\leq 0$ and denominator $> 0$, so it is $\leq 0$; since the two-sided limit exists, it equals the right-hand limit, and non-strict inequalities survive limits, so $f'(c) \leq 0$. For $x$ slightly to the left, numerator $\leq 0$ and denominator $< 0$ give a quotient $\geq 0$, so $f'(c) \geq 0$. Hence $f'(c) = 0$. Replace $f$ by $-f$ for a minimum.
\end{proof}

Two things to internalise. First, \emph{strict} inequalities do not survive limits---$1/n > 0$ for all $n$, and the limit is not positive---but non-strict ones do; the proof would collapse if it needed the strict version. Second, differentiability is used twice, once per side, which is why the lemma is false at endpoints. The function $f(x) = x$ on $[0,1]$ has a maximum at $1$ and no vanishing derivative anywhere.

\begin{theorem}[Rolle]\label{thm:rolle}
If $f$ is continuous on $[a,b]$, differentiable on $(a,b)$, and $f(a) = f(b)$, then $f'(\xi) = 0$ for some $\xi \in (a,b)$.
\end{theorem}

\begin{proof}
Theorem~\ref{thm:evt} supplies a maximiser $q$ and a minimiser $p$ in $[a,b]$. If both are endpoints, then $\max f = \min f = f(a)$, so $f$ is constant and every interior point works. Otherwise one of them lies in $(a,b)$, and Lemma~\ref{lem:fermat} applies to it.
\end{proof}

Note the asymmetry of hypotheses: continuity on the closed interval, differentiability only on the open one. This is not fastidiousness. It is what allows the theorem to apply to $\sqrt{1 - x^{2}}$ on $[-1,1]$, whose derivative is infinite at both ends.

\begin{theorem}[Mean Value Theorem]\label{thm:mvt}
If $f$ is continuous on $[a,b]$ and differentiable on $(a,b)$, then
\[
f'(\xi) = \frac{f(b) - f(a)}{b - a}
\]
for some $\xi \in (a,b)$.
\end{theorem}

\begin{proof}
Subtract off the secant. Let
\[
h(x) = f(x) - f(a) - \frac{f(b) - f(a)}{b - a}\,(x - a).
\]
Then $h$ inherits continuity on $[a,b]$ and differentiability on $(a,b)$ from $f$ and an affine function,
\[
h'(x) = f'(x) - \frac{f(b) - f(a)}{b - a},
\]
and $h(a) = h(b) = 0$. Theorem~\ref{thm:rolle} produces $\xi$ with $h'(\xi) = 0$.
\end{proof}

The Mean Value Theorem is Rolle in a tilted coordinate frame. Nothing is added; the graph is merely sheared until the chord is horizontal. Once you see that, the two theorems stop being separate items to memorise.

\section{What the structure was}

The dependency order is a partial one, not a chain:
\[
\text{Completeness} \longrightarrow \{\,\text{IVT},\ \text{EVT}\,\},
\qquad
\text{EVT} + \text{Lemma \ref{lem:fermat}} \longrightarrow \text{Rolle} \longrightarrow \text{MVT}.
\]
The Intermediate and Extreme Value Theorems are independent; neither is used in proving the other, and the former is used nowhere in the derivation of the Mean Value Theorem. Textbooks that invoke it inside a Mean Value Theorem proof are importing a hypothesis they do not need. The two do combine, though, to something stronger than either: a continuous image of $[a,b]$ is exactly $[\min f, \max f]$---a compact interval maps onto a compact interval. That statement is the real theorem of which the other two are the halves.

There are shorter routes. The Extreme Value Theorem falls out of Bolzano--Weierstrass in four lines: take $x_{n}$ with $f(x_{n}) \to \sup f$, extract a convergent subsequence, use sequential continuity. The proofs here avoid that deliberately, because Bolzano--Weierstrass is itself a repackaging of completeness, and the point was to keep the axiom visible rather than to be brief. Later, when you have compactness as a concept, all of Section~4 becomes one sentence about continuous images of compact sets. That compression is what the abstraction buys you, and it is only meaningful if you have first paid for it in supremum arguments.

Finally, the fact worth carrying: over an ordered field, each of these three theorems is \emph{equivalent} to the least-upper-bound property. Assume the Intermediate Value Theorem and you can derive completeness. They are not consequences of the real numbers; they are the real numbers, stated in the language of functions. When a proof in this essay felt like it was doing nothing but manipulating $\sup$, that was not a stylistic failure. That was the content.

\end{document}
